\section{Interpretación y decisiones}
El diagnóstico integral del caso muestra una liquidez razonable, una estructura de financiamiento con participación importante de terceros y una rentabilidad positiva. Sin embargo, el análisis no debe concluir que la empresa se encuentra automáticamente en buena situación, porque faltan comparativos sectoriales, metas internas y periodos anteriores.

\begin{table}[h]
\centering
\begin{tabular}{p{0.22\textwidth}p{0.33\textwidth}p{0.33\textwidth}}
\toprule
Área & Interpretación & Decisión \\
\midrule
Liquidez & Razón corriente 1.80 y prueba ácida 1.20. & Vigilar cobranza e inventarios. \\
Gestión & Requiere promedios para evaluar ciclo operativo. & Medir cobranza e inventario mensualmente. \\
Solvencia & 47.50 \char37{} de activos financiados por terceros. & Revisar costo financiero y vencimientos. \\
Rentabilidad & Margen neto 8.46 \char37{}. & Analizar margen por producto y gastos. \\
\bottomrule
\end{tabular}
\end{table}

Las decisiones recomendadas son fortalecer cobranza, controlar inventarios, revisar costo financiero, evaluar margen por producto, proyectar flujo de efectivo y monitorear indicadores mensualmente.
